% ===== tr-03 块2：第 3 章 pp.26--28（原文页 26--28）=====
% [原文 p.26]
对偶相交复形是不是正确的带奇点仿射流形？下面的定理为此提供了证据，并给出了这一构造与 SYZ 猜想之间的联系。

\begin{theorem}{3.7.7}
设 $\mathcal X\to D$ 是环面退化，其对偶相交复形为 $(B,\mathcal P)$。则存在开集 $U\subseteq B$，使得 $B\setminus U$ 收缩到 $B$ 的判别轨迹 $\Gamma$ 上，并且 $\mathcal X_t$ 含有一个开子集 $\mathcal U_t$，它作为复流形同构于 $X_\epsilon(U)$ 的某个扭转的一个小形变，其中 $\epsilon=O(-1/\ln|t|)$。
% [Tutor Note: 这里的开子集记作花体 $\mathcal U_t$，不要与定理中开集 $U\subseteq B$ 的记号混淆；$\epsilon=O(-1/\ln|t|)$ 中的 ln 是自然对数，文字层常把 $|t|$ 抽成 t]
\end{theorem}

这里我们不会精确说明所谓小形变是什么意思；所谓扭转，指的是用 $B$-场对 $X_\epsilon(U)$ 的复结构作扭转。精确得多的陈述见 [20]；上面的陈述只是想给出对真实情形的感觉。其证明连同精确得多的陈述，最终将出现在 [21] 中。

若 $\mathcal X\to D$ 是\emph{极化}环面退化，即 $\mathcal X$ 上存在相对丰沛的线丛 $\mathcal L$，则我们可以如下构造另一个带奇点的仿射流形与多面体分解 $(\check B,\check{\mathcal P})$，称之为\emph{相交复形}。

对 $\mathcal X_0$ 的每个不可约分支 $X_i$，$\mathcal L|_{X_i}$ 是环面簇上的一个丰沛线丛。以 $\sigma_i\subseteq N_\R$ 记该线丛的牛顿多面体。于是 $\mathcal X_0$ 中容于 $X_i$ 的分层与 $\sigma_i$ 的面之间存在一一的、保包含关系的对应。接着我们可以按显然的方式把各个 $\sigma_i$ 粘合起来：若 $Y$ 是 $\mathcal X_0$ 的余维 1 分层，则它含于两个不可约分支 $X_i$ 与 $X_j$ 中，并分别确定 $\sigma_i$ 与 $\sigma_j$ 的面。这些面是仿射同构的，因为它们同为 $\mathcal L|_Y$ 的牛顿多面体，于是我们可以按典范方式把它们等同起来。这样便得到一个带多面体分解 $\check{\mathcal P}$ 的拓扑空间 $\check B$。我们按与前面类似的方式给它一个带奇点的仿射结构。同样地，设 $\Gamma$ 是 $\operatorname{Bar}(\check{\mathcal P})$ 中那些不含 $\check{\mathcal P}$ 的顶点、或不与 $\check{\mathcal P}$ 的某个极大胞腔的内部相交的单形之并。置 $\check B_0:=\check B\setminus\Gamma$，它同样有开覆盖
\[
\{W_\sigma\mid\sigma\in\check{\mathcal P}\ \text{极大}\}\cup\{W_v\mid v\in\check{\mathcal P}\ \text{顶点}\}.
\]
照例，由于 $W_\sigma$ 是 $\sigma$ 的内部，它自带一个典范仿射结构。另一方面，$\check{\mathcal P}$ 的顶点 $v$ 对应 $\mathcal X_0$ 的一个零维分层 $x$，而联系于 $x$ 的是多面体 $\sigma_x\subseteq M_\R$。设 $\check\Sigma_x$ 是 $\sigma_x$ 在 $N_\R$ 中的正规扇。于是 $\check\Sigma_x$ 中的锥与 $\mathcal X_0$ 中含 $x$ 的分层之间存在一一的、保包含关系的对应。这一对应允许我们定义一个卡
\[
\check\psi_v:W_v\to N_\R,
\]
它把 $W_v\cap\check\sigma$（对任何以 $v$ 为顶点、属于 $\check{\mathcal P}$ 的 $\check\sigma$）整仿射线性地映入 $\check\Sigma_x$ 的相应锥中。这就显然在 $\check B_0$ 上给出了整仿射结构，从而定义了相交复形 $(\check B,\check{\mathcal P})$。
% [Tutor Note: 全文的 check（ˇ）一律表示"对偶"侧的记号，写在符号整体之上：$\check B$、$\check{\mathcal P}$、$\check\sigma$、$\check\Sigma_x$；它不是 tilde，也不要写成 $\mathcal P^\vee$ 之类]

与定理 3.7.7 类似，我们期望

\begin{conjecture}{3.7.8}
设 $\mathcal X\to D$ 是极化环面退化，其相交复形为 $(\check B,\check{\mathcal P})$。设 $\omega_t$ 是 $\mathcal X_t$ 上代表该极化的第一 Chern 类的 Kähler 形式。则存在开集 $\check U\subseteq\check B$，使得 $\check B\setminus\check U$ 收缩到 $\check B$ 的判别轨迹 $\Gamma$ 上，并且对任意 $t$，$\mathcal X_t$ 都是 $\check X(\check U)$ 的辛紧化。
\end{conjecture}

% [原文 p.27]
我并不认为这特别困难：它应当可以用 Ruan [53] 的技巧处理，但尚未有人完成。

通过引入多值分段线性函数，可以把相交复形与对偶相交复形之间的关系刻画得更精确；这与定义 3.2.3 中的多值凸函数相类比：

\begin{definition}{3.7.9}
设 $B$ 是带奇点的仿射流形，带多面体分解 $\mathcal P$。则 $B$ 上的一个多值分段线性函数 $\varphi$ 是指开覆盖 $\{(U_i,\varphi_i)\}$ 上的一族连续函数，使得 $\varphi_i$ 在 $\mathcal P$ 的每个与 $U_i$ 相交的胞腔上为仿射线性，并且在 $U_i\cap U_j$ 上 $\varphi_i-\varphi_j$ 为仿射线性。此外，对任意 $\sigma\in\mathcal P$，在每一点 $x\in U_i\cap\operatorname{Int}(\sigma)$ 的某个邻域内，存在仿射线性函数 $\psi$，使得 $\varphi_i-\psi$ 在 $\sigma$ 上为零。
\end{definition}

为解释最后这一条件、并澄清 $B$ 上的附加结构，我们来考察当 $(B,\mathcal P)$ 是对偶相交复形时，$B$ 的多面体分解 $\mathcal P$ 所具有的一个性质。

考虑任意 $p$ 维胞腔 $\sigma\in\mathcal P$。它对应一个 $n-p$ 维分层 $X_\sigma\subseteq\mathcal X_0$，因而是一个由 $\R^{n-p}$ 中的扇 $\Sigma_\sigma$ 定义的环面簇。现在，对 $\sigma$ 的任意顶点 $v$，$X_\sigma$ 是 $\mathcal X_0$ 的不可约分支 $X_v$ 的一个环面分层。于是 $\Sigma_\sigma$ 可作为 $\Sigma_v$ 的\emph{商扇}得到。换言之，存在 $\Sigma_v$ 中对应 $\sigma$ 的 $p$ 维锥 $K_\sigma$，使得
\[
\Sigma_\sigma=\Sigma_v(K_\sigma):=\{(K+\R K_\sigma)/\R K_\sigma\mid K\in\Sigma_v,\ K\supseteq K_\sigma\}.
\]
特别地，存在 $\operatorname{Int}(K_\sigma)$ 的一个开邻域 $U_{v,\sigma}$ 与一个整线性映射 $S_{v,\sigma}:U_{v,\sigma}\to\R^{n-p}$，使得
\[
\{S_{v,\sigma}^{-1}(K)\mid K\in\Sigma_\sigma\}=\{U_{v,\sigma}\cap K\mid K\in\Sigma_v,\ K\supseteq K_\sigma\}.
\]
设 $U_\sigma$ 是 $B$ 中 $\operatorname{Int}(\sigma)$ 的一个小开邻域；若取得足够小，映射 $S_{v,\sigma}$ 可看作定义在 $U_\sigma$ 的开子集上，并能拼贴起来给出一个整仿射浸没 $S_\sigma:U_\sigma\to\R^{n-p}$，其中 $U_\sigma=\bigcup_{v\in\sigma}U_{v,\sigma}$ 是 $\operatorname{Int}(\sigma)$ 的一个开邻域。这个映射具有性质
\[
\{S_\sigma^{-1}(K)\mid K\in\Sigma_\sigma\}=\{U_\sigma\cap\tau\mid \tau\supseteq\sigma,\ \tau\in\mathcal P\}.
\]
一般地，我们称一个多面体分解是\emph{环面的}，如果对所有的 $\sigma\in\mathcal P$ 总存在这样从 $\operatorname{Int}(\sigma)$ 的邻域 $U_\sigma$ 出发的整仿射线性映射 $S_\sigma:U_\sigma\to\R^{n-p}$，以及 $\R^p$ 中具有上述性质的扇 $\Sigma_\sigma$。
% [Tutor Note: 上一句原文作 a fan $\Sigma_\sigma$ in $\R^p$，但 $\Sigma_\sigma$ 本节开头明确说位于 $\R^{n-p}$，且 $S_\sigma$ 的靶空间是 $\R^{n-p}$，此处应为 $\R^{n-p}$，属原文笔误（550 dpi 裁图核对确印作 $\R^p$）；正文照原文排]
（关于环面多面体分解的一个或许过于精细的定义，见 [22] 的定义 1.22。那里的定义由于允许胞腔自相交、等价地允许 $\mathcal X_0$ 的不可约分支非正规，而变得复杂。）我们可以把扇 $\Sigma_\sigma$ 看作 $\mathcal P$ 在 $\sigma$ 内部某点处横截于 $\sigma$ 的扇结构。对偶相交复形的要点在于：这个扇结构由 $X_\sigma$ 决定，且与 $\sigma$ 内部所取点的选择无关。

现在回到分段线性函数。设给定极化环面退化 $\mathcal X\to D$。事实上，我们可以如下在对偶相交复形 $(B,\mathcal P)$ 上得到一个分段线性函数 $\varphi$。限制到任意环面分层 $X_\sigma$ 上，$\mathcal L|_{X_\sigma}$ 完全由 $\Sigma_\sigma$ 上的一个整分段线性函数 $\bar\varphi_\sigma$ 决定，且该函数在相差一个线性函数的意义下良定。把这个分段线性函数经 $S_\sigma$ 拉回到 $U_\sigma$，便得到一族分段线性函数 $\{(U_\sigma,\varphi_\sigma)\mid\sigma\in\mathcal P\}$。对 $\tau\subseteq\sigma$ 有 $(\mathcal L|_{X_\tau})|_{X_\sigma}=\mathcal L|_{X_\sigma}$ 这一事实蕴涵：在重叠处 $\varphi_\sigma$ 与 $\varphi_\tau$ 至多相差一个线性函数。
% [原文 p.28]
于是 $\{(U_\sigma,\varphi_\sigma)\}$ 定义了多值分段线性函数。多值分段线性函数定义中的最后一个条件于是反映出：该函数在 $\sigma$ 的邻域内必须是经 $S_\sigma$ 拉回的函数。事实上，给定 $(B,\mathcal P)$ 上任一多值分段线性函数 $\varphi$（其中 $\mathcal P$ 是 $B$ 的环面多面体分解），$\varphi$ 都由 $\Sigma_\sigma$（$\sigma\in\mathcal P$）上的函数 $\bar\varphi_\sigma$ 经 $S_\sigma$ 拉回完全决定。

若 $\mathcal L$ 丰沛，则由 $\mathcal L|_{X_\sigma}$ 决定的分段线性函数是严格凸的。于是我们说一个多值分段线性函数是\emph{严格凸的}，如果每个 $\sigma\in\mathcal P$ 的 $\bar\varphi_\sigma$ 都严格凸。

现在抽象地给定三元组 $(B,\mathcal P,\varphi)$，其中 $B$ 是带奇点的整仿射流形、配有环面多面体分解 $\mathcal P$，而 $\varphi$ 是 $B$ 上的严格凸多值分段线性函数。则我们如下构造 $(B,\mathcal P,\varphi)$ 的\emph{离散 Legendre 变换} $(\check B,\check{\mathcal P},\check\varphi)$。

$\check B$ 将通过粘合牛顿多面体来构造。对 $\mathcal P$ 的顶点 $v$，若把扇 $\Sigma_v$ 看作位于 $M_\R$ 中，则 $\bar\varphi_v$ 的牛顿多面体为
\[
\check v=\{x\in N_\R\mid\langle x,y\rangle\ge-\bar\varphi_v(y)\ \ \forall y\in M_\R\}.
\]
$\check v$ 的面与 $\mathcal P$ 中含 $v$ 的胞腔之间存在一一的、序反转的对应。此外，若 $\sigma$ 是 $\mathcal P$ 中含两个顶点 $v$ 与 $v'$ 的最小胞腔，则 $\check v$ 与 $\check v'$ 的对应面是整仿射同构的，因为它们都同构于 $\bar\varphi_\sigma$ 的牛顿多面体。于是我们可以沿这个公共面把 $\check v$ 与 $\check v'$ 粘合起来。作完所有这些等同之后，我们得到一个胞腔复形 $(\check B,\check{\mathcal P})$，它其实正是 $(B,\mathcal P)$ 的对偶胞腔复形。当然，我们还有一些附加信息，即 $\check{\mathcal P}$ 的每个极大胞腔内部上的一个仿射结构。为给 $\check B$ 一个带奇点的整仿射结构，照例进行：使用这个仿射结构，并把 $\check{\mathcal P}$ 的每个顶点的邻域与 $\mathcal P$ 中相应极大胞腔的正规扇等同起来。
% [Tutor Note: $\check v$ 的定义里 $-\bar\varphi_v(y)$ 的负号极易漏抄；另外 $\Sigma_v$ 位于 $M_\R$，而牛顿多面体 $\check v$ 位于对偶空间 $N_\R$，两者不要写反]

最后，函数 $\varphi$ 在 $(\check B,\check{\mathcal P})$ 上有一个离散 Legendre 变换 $\check\varphi$。我们别无选择，只能这样定义：在顶点 $\check\sigma\in\check{\mathcal P}$ 的邻域内（其中 $\check\sigma$ 对偶于极大胞腔 $\sigma\in\mathcal P$），取 $\check\varphi$ 为牛顿多面体是 $\sigma$ 的分段线性函数，即
\[
\check\varphi_{\check\sigma}(y)=-\inf\{\langle y,x\rangle\mid x\in\sigma\subseteq M_\R\}.
\]
这给出 $(\check B,\check{\mathcal P},\check\varphi)$，即 $(B,\mathcal P,\varphi)$ 的离散 Legendre 变换。若 $B$ 是 $\R^n$，则它与离散 Legendre 变换的经典概念一致。离散 Legendre 变换有几个相关的性质：
\begin{itemize}
\item $(\check B,\check{\mathcal P},\check\varphi)$ 的离散 Legendre 变换是 $(B,\mathcal P,\varphi)$。
\item 若把底拓扑空间 $B$ 与 $\check B$ 通过"它们是互相对偶的胞腔复形的底空间"而等同起来，则 $\Lambda_{B_0}\cong\check\Lambda_{\check B_0}$ 且 $\check\Lambda_{B_0}\cong\Lambda_{\check B_0}$，其中下标指明用哪个仿射结构来定义 $\Lambda$ 或 $\check\Lambda$。
\end{itemize}

这有望说明：离散 Legendre 变换是第 2 节中 Legendre 变换所提供的对偶性的合适替代。

最后，它引向我们或可称之为\emph{代数几何的 SYZ 程序}的东西。与第 5 节所建议的程序相类比，我们依以下步骤进行：
